\section{DeepSeek 与近期 scaling recipes}

MiniCPM 展示了依赖 muP 的路线，接下来 DeepSeek 提供另一种工程选择：不要求所有超参理论上 scale-invariant，而是直接通过小规模 runs 估计 optimal batch、learning rate 与 IsoFLOP model/data 配比。比较两者能看出，scaling recipe 的关键是可验证预测，而不是统一采用某个参数化流派。


\singleslide{slides-images/slide-019.jpg}{DeepSeek：另一套公开 scaling analysis，用 7B/67B 高性能模型检验 MiniCPM 规律。}{19}
跨团队复现能区分普遍规律与单一实现细节，尤其是 LR、batch 与 data/model allocation。
\begin{knowledgebox}{讲义补充：源 Slide 19 的 DeepSeek 对照}
DeepSeek 与 MiniCPM 的区别在于它不依赖 muP，而是直接估计 batch/LR scaling，并使用 WSD-style learning rate 和 IsoFLOP 分析。它代表另一条路线：用足够多小规模实验证明超参趋势，而不是依赖参数化理论稳定性。
\end{knowledgebox}

\begin{figure}[H]
\centering
\includegraphics[width=0.86\textwidth]{slide-020.jpg}
\caption{Slide 20：DeepSeek scaling strategy：不用 muP，直接估计 optimal batch 和 LR。}
\figsource{CS336 Spring 2026 Lecture 11 官方幻灯片。}
\end{figure}

\begin{warningbox}{读图：Slide 20 的风险}
不用 muP 意味着最优 LR 可能随 scale 漂移，因此必须更认真地做 LR/batch sweep。直接经验拟合很实用，但外推风险也更依赖实验覆盖和函数形式选择。
\end{warningbox}

\begin{figure}[H]
\centering
\includegraphics[width=0.86\textwidth]{slide-021.jpg}
\caption{Slide 21：DeepSeek 的 learning rate scaling，用小规模 runs 收集 near-optimal 模型；LR fit 看起来有些可疑。}
\figsource{CS336 Spring 2026 Lecture 11 官方幻灯片。}
\end{figure}

\begin{knowledgebox}{读图：Slide 21 的“questionable fit”怎么看}
Near-optimal 定义为距离最小 loss 0.25\% 以内，这能缓解噪声，但 LR 的最优区间可能很平。若点分布宽、拟合斜率不稳，就不能过度相信外推公式。公开论文里的 scaling fit 也需要质疑。
\end{knowledgebox}

\begin{figure}[H]
\centering
\includegraphics[width=0.86\textwidth]{slide-022.jpg}
\caption{Slide 22：DeepSeek 使用 WSD-style learning rate，fast warmup 加两次 10\% decay，性能接近 cosine。}
\figsource{CS336 Spring 2026 Lecture 11 官方幻灯片。}
\end{figure}

\begin{knowledgebox}{读图：Slide 22 的 WSD 共性}
MiniCPM 和 DeepSeek 都使用 WSD 类思想，说明 WSD 的价值不仅是曲线形状，而是可复用训练轨迹、支持 scaling 分析。它让同一段 stable training 可以接多个 decay endings，减少重复从头训练。
\end{knowledgebox}

\begin{figure}[H]
\centering
\includegraphics[width=0.86\textwidth]{slide-023.jpg}
\caption{Slide 23：DeepSeek 的 data-size tradeoff 使用 Chinchilla method 2，即 IsoFLOP 风格分析。}
\figsource{CS336 Spring 2026 Lecture 11 官方幻灯片。}
\end{figure}

\begin{importantbox}{读图：Slide 23 的 IsoFLOP 逻辑}
固定 FLOP budget，扫 model size 和 data size，取 loss 最小点。这样得到每个 compute 下的 optimal model size。它比直接 joint fit 更直观，也更容易发现某些点是否异常。
\end{importantbox}

\begin{figure}[H]
\centering
\includegraphics[width=0.86\textwidth]{slide-024.jpg}
\caption{Slide 24：DeepSeek 的 fitted scaling models 通常能准确预测 final model losses。}
\figsource{CS336 Spring 2026 Lecture 11 官方幻灯片。}
\end{figure}

\begin{knowledgebox}{读图：Slide 24 是 scaling recipe 的验收}
Scaling law 最终要看外推预测是否命中大模型 final loss。若小规模 fitted curve 能预测 7B/67B 等最终 loss，说明超参、数据、训练 recipe 在外推区间内足够稳定。反之，漂亮的小规模曲线也可能没有工程价值。
\end{knowledgebox}

\subsection{更多公开模型 scaling 片段}

为了判断 DeepSeek 式 recipe 是否只是个例，本节横向查看 Qwen、Kimi、Hunyuan、LLaMA 与 MiniMax 等公开材料。不同团队拟合的对象包括 dense/MoE 参数量、sparsity、batch、learning rate 与 downstream metric；阅读重点是实验自变量、active parameter 定义和预测是否经过最终规模验证。


\singleslide{slides-images/slide-025.jpg}{Qwen 2.5/3 的 scaling evidence：batch 与 hyperparameter fits 仍是近期 recipe 的常见组成。}{25}
这页说明 scaling practice 没有在 Chinchilla 后停止，而是扩展到优化超参与不同模型族。
\begin{knowledgebox}{讲义补充：源 Slide 25 的 Qwen 证据}
Qwen 的公开信息说明 batch 和 LR scaling 已经成为主流模型训练报告的重要组成。不同团队可能用不同函数形式，但共同目标都是减少大模型调参风险。
\end{knowledgebox}


\singleslide{slides-images/slide-026.jpg}{Kimi K2：用 sparsity scaling law 选择 MoE 的稀疏程度。}{26}
MoE 不只决定总参数，还要选择 activated experts/parameters；scaling law 可比较质量收益与路由、通信和存储成本。
\begin{knowledgebox}{讲义补充：源 Slide 26 把 scaling 扩展到 sparsity}
Scaling 不只决定模型大小和 token 数，也可以决定 sparse/MoE 模型里 active parameters、expert 数、稀疏比例。Kimi K2 这类模型说明“参数价值”在稀疏模型里更复杂，必须把 active compute 和 total parameters 分开。
\end{knowledgebox}


\singleslide{slides-images/slide-027.jpg}{Hunyuan MoE IsoFLOPS：拟合 data 与 active-parameter 的最优约 96:1 比例。}{27}
对 MoE 应使用每 token active parameters 解释 compute，同时另记 total parameters 的 memory/communication 成本。
\begin{knowledgebox}{讲义补充：源 Slide 27 的 active param ratio}
MoE 的 total parameters 远大于每 token active parameters，因此 compute-optimal ratio 应以 active param 或 activated FLOPs 理解。96:1 这样的比例提醒我们，Chinchilla dense 模型口径不能直接套到 MoE。
\end{knowledgebox}

\begin{figure}[H]
\centering
\includegraphics[width=0.86\textwidth]{slide-028.jpg}
\caption{Slide 28：LLaMA 3 scaling laws，包括 IsoFLOPS-style scaling、39:1 ratio 和 compute-to-downstream scaling。}
\figsource{CS336 Spring 2026 Lecture 11 官方幻灯片。}
\end{figure}

\begin{knowledgebox}{读图：Slide 28 的两个层次}
预训练 loss 的 IsoFLOPS 只是第一层；compute-to-downstream scaling 关注下游任务如何随 compute 改善。大模型训练真正关心的是产品能力和 benchmark，而不仅是 cross entropy。因此 downstream scaling 是更难也更有价值的外推。
\end{knowledgebox}


\singleslide{slides-images/slide-029.jpg}{MiniMax-01：把 architecture choices 也纳入 scaling law，并结合 Chinchilla Method 1。}{29}
当 attention、depth、width 或 context 机制改变时，不能假设旧 exponent 不变；需要模型族级对照实验。
\begin{knowledgebox}{讲义补充：源 Slide 29 的 architecture scaling}
Architecture scaling laws 用来比较不同结构随规模变化的 loss 或能力曲线。MiniMax-01 的例子说明，架构选择也可以进入 scaling law pipeline，而不只是凭最大模型结果判断。
\end{knowledgebox}

\begin{figure}[H]
\centering
\includegraphics[width=0.86\textwidth]{slide-030.jpg}
\caption{Slide 30：近期 scaling recipe 总结，以 DeepSeek 为代表：假设大部分 Transformer hypers scale-invariant，重点拟合 batch/LR 和 model sizing。}
\figsource{CS336 Spring 2026 Lecture 11 官方幻灯片。}
\end{figure}

\begin{importantbox}{读图：Slide 30 的 recipe 可执行版本}
实践中常先固定大多数 architecture hypers，只做 batch/LR scaling，再用 IsoFLOP 选择模型大小和 token 数。这样把指数级超参搜索压缩成少数关键轴。代价是：如果固定的 architecture hypers 实际随 scale 变化，recipe 会漏掉更优方案。
\end{importantbox}

\subsection{本章小结}

MiniCPM 和 DeepSeek 展示了两种 scaling 实践：一类用 muP 稳定迁移，另一类用直接 empirical sweep 拟合 batch/LR 和 IsoFLOP。近期模型则说明 scaling law 的对象正在扩展到 MoE sparsity、architecture、downstream 能力和 active parameters。

